% ============================================================
%  chapters/ch-logic-boolean-geometry.tex
%  Part VIII: Mathematical Foundations — Logic and Boolean Geometry
% ============================================================

\chapter{Logic as Geometry: The Boolean Landscape}

\chapterepigraph{%
  Logic is not the science of valid inference.\\
  That is merely its application.\\
  Logic is the geometry of truth.
}{Flyxion, \textit{Semantic Infrastructure}}

\sidenote{App.\,D and\\App.\,E give the\\complete formal\\treatment;\\ connects to\\Boolean lattice\\diagram (Fig.\,\ref{fig:hasse})}

Part~VIII occupies an unusual position in the monograph. The material
here — Boolean lattices, hypercube decompositions, the geometry of
logical operators — was historically developed \emph{before} the
frameworks of Parts~I–VII. But it is placed here because it can now
be understood in its proper context: as a special case of the
accessibility geometry developed in Part~IV, specialized to the
discrete two-valued domain.

Logical operators are not isolated symbols. They are points in a
geometric space. Understanding them geometrically reveals structure
that the symbolic presentation obscures.

\section{The Curriculum Question}

Before developing the geometry, we should address a pedagogical
question that this chapter raises directly. The standard mathematics
curriculum is \emph{algebra-first}: children encounter numbers,
operations, equations, and functions before they encounter logic,
set theory, or geometry as separate disciplines. Logic, if it appears
at all, arrives late — as a chapter in a discrete mathematics course,
or as the preamble to a proof-writing course.

What would a \emph{logic-first} curriculum look like? And what kind
of mathematical thinking would it produce?

\begin{conceptaside}{Seven Mathematical Schools}
  Consider seven possible entry points into mathematics, each
  corresponding to a different primary mathematical object:
  \begin{itemize}
    \item \textbf{Geometry school:} space, shape, symmetry, transformation.
    \item \textbf{Algebra school:} symbol, equation, structure, abstraction.
    \item \textbf{Trigonometry school:} cycle, wave, angle, periodicity.
    \item \textbf{Calculus school:} change, rate, accumulation, limit.
    \item \textbf{Statistics school:} uncertainty, evidence, inference,
          variation.
    \item \textbf{Stoichiometry school:} conservation, transformation,
          balance, reaction.
    \item \textbf{Logic school:} truth, implication, consistency, proof.
  \end{itemize}
  Each school produces a different cognitive orientation. The logic-first
  student learns to ask: \emph{what follows necessarily from what?}
  The geometry-first student asks: \emph{what is the shape of this?}
  The statistics-first student asks: \emph{what is the evidence for this?}
  These orientations are complementary — the richest mathematical minds
  move fluidly between them — but they are genuinely different.
\end{conceptaside}

\section{The Space of Binary Connectives}

The starting point of Boolean geometry is the observation that there
are exactly 16 binary connectives — 16 Boolean functions of two variables.
These 16 functions form a geometric object: the four-dimensional hypercube
$Q_4$.

\begin{definition}{Binary Connective Space}{binary-connective}
  The \emph{binary connective space} is $B_2 = \{0,1\}^4$, with each
  vertex corresponding to one binary connective, labeled by its output
  column over the four input states $TT, TF, FT, FF$.
\end{definition}

The familiar operators sit at specific vertices:
\begin{align*}
  \top &= 1111 & \bot &= 0000 \\
  \mathrm{AND} &= 0001 & \mathrm{OR} &= 0111 \\
  \mathrm{NAND} &= 1110 & \mathrm{NOR} &= 1000 \\
  \mathrm{XOR} &= 0110 & \mathrm{XNOR} &= 1001
\end{align*}

The Hamming distance between vertices is their \emph{logical distance}:
how many truth-value assignments they disagree on.

\begin{figure}[ht]
  \centering
  \begin{tikzpicture}[scale=0.85]
    \tikzset{
  boolnode/.style={circle, draw=RSVPdeep, fill=white, minimum size=8mm,
    font=\tiny\ttfamily, line width=0.6pt},
  namednode/.style={circle, draw=RSVPdeep, fill=RSVPlight, minimum size=8mm,
    font=\tiny\bfseries\ttfamily, line width=1.2pt},
  topbot/.style={circle, fill=RSVPdeep, text=white, minimum size=8mm,
    font=\tiny\bfseries\ttfamily, line width=1.2pt},
  edge/.style={RSVPmid!50, thin},
  namedge/.style={RSVPaccent, thick, opacity=0.8},
  lbl/.style={font=\scriptsize\itshape, RSVPmid},
}
% ============================================================
%  diagrams/boolean-hasse-lattice.tex
%  The 16 binary Boolean functions as a Hasse diagram
%  Nodes labeled by 4-bit truth table output
%  Identified named connectives highlighted
% ============================================================

  node distance=1.3cm and 1.1cm,
  boolnode/.style={
    circle, draw=RSVPdeep, fill=white,
    minimum size=8mm, font=\tiny\ttfamily,
    line width=0.6pt
  },
  namednode/.style={
    boolnode, fill=RSVPlight, draw=RSVPdeep, line width=1.2pt,
    font=\tiny\bfseries\ttfamily
  },
  topbot/.style={
    boolnode, fill=RSVPdeep, text=white, line width=1.2pt,
    font=\tiny\bfseries\ttfamily
  },
  edge/.style={RSVPmid!50, thin},
  namedge/.style={RSVPaccent, thick, opacity=0.8},
  lbl/.style={font=\scriptsize\itshape, RSVPmid}
% ── Level 0: Contradiction (bottom) ──────────────────────────
\node[topbot] (n0) at (0,0) {0000};
\node[lbl, below=2pt of n0] {$\bot$};

% ── Level 1: four 1-element functions ─────────────────────────
\node[boolnode] (n1) at (-4.5, 1.6) {0001};   % NOR
\node[boolnode] (n2) at (-1.5, 1.6) {0010};
\node[boolnode] (n4) at ( 1.5, 1.6) {0100};
\node[boolnode] (n8) at ( 4.5, 1.6) {1000};   % AND
\node[lbl, left=2pt of n1] {NOR};
\node[lbl, right=2pt of n8] {AND};

% ── Level 2: six 2-element functions ─────────────────────────
\node[namednode] (n3) at (-5, 3.4) {0011};    % ¬A
\node[boolnode]  (n5) at (-3, 3.4) {0101};
\node[namednode] (n6) at (-1, 3.4) {0110};    % XOR
\node[boolnode]  (n9) at ( 1, 3.4) {1001};    % XNOR
\node[boolnode]  (nA) at ( 3, 3.4) {1010};
\node[namednode] (nC) at ( 5, 3.4) {1100};    % ¬B
\node[lbl, left=2pt of n3] {$\neg A$};
\node[lbl, below right=0pt and 1pt of n6] {XOR};
\node[lbl, right=2pt of nC] {$\neg B$};

% ── Level 3: six 3-element functions ─────────────────────────
\node[namednode] (n7) at (-4.5, 5.2) {0111};  % NAND
\node[boolnode]  (nB) at (-1.5, 5.2) {1011};
\node[namednode] (nD) at ( 1.5, 5.2) {1101};
\node[namednode] (nE) at ( 4.5, 5.2) {1110};  % OR
\node[lbl, left=2pt of n7] {NAND};
\node[lbl, right=2pt of nE] {OR};

% ── Level 4: Tautology (top) ─────────────────────────────────
\node[topbot] (nF) at (0, 6.8) {1111};
\node[lbl, above=2pt of nF] {$\top$};

% ── Edges (level 0 → 1) ───────────────────────────────────────
\foreach \a/\b in {n0/n1, n0/n2, n0/n4, n0/n8}{
  \draw[edge] (\a) -- (\b);
}

% ── Edges (level 1 → 2) ───────────────────────────────────────
\draw[edge] (n1) -- (n3); \draw[edge] (n1) -- (n5);
\draw[edge] (n2) -- (n3); \draw[edge] (n2) -- (n6);
\draw[edge] (n4) -- (n5); \draw[edge] (n4) -- (nC);
\draw[edge] (n8) -- (n9); \draw[edge] (n8) -- (nC);
\draw[edge] (n2) -- (nA); \draw[edge] (n8) -- (nA);
\draw[edge] (n1) -- (n9); \draw[edge] (n4) -- (n6);

% ── Edges (level 2 → 3) ───────────────────────────────────────
\draw[edge] (n3) -- (n7); \draw[edge] (n5) -- (n7);
\draw[edge] (n6) -- (n7); \draw[edge] (n6) -- (nE);
\draw[edge] (n9) -- (nB); \draw[edge] (nA) -- (nB);
\draw[edge] (n3) -- (nB); \draw[edge] (nC) -- (nE);
\draw[edge] (nA) -- (nD); \draw[edge] (n5) -- (nD);
\draw[edge] (n9) -- (nD); \draw[edge] (nC) -- (nB);

% ── Edges (level 3 → 4) ───────────────────────────────────────
\foreach \a in {n7, nB, nD, nE}{
  \draw[edge] (\a) -- (nF);
}

% ── Highlighted path: ⊥ → NOR → ¬A → NAND → ⊤ ──────────────
\draw[namedge] (n0) -- (n1) -- (n3) -- (n7) -- (nF);

% ── Highlighted path: ⊥ → AND → ¬B → OR → ⊤ ─────────────────
\draw[RSVPpurple, thick, opacity=0.7] (n0) -- (n8) -- (nC) -- (nE) -- (nF);

% ── Hypercube level labels ────────────────────────────────────
\node[lbl, right=5.8cm of n0, anchor=west] {Level 0: Contradiction};
\node[lbl, right=5.8cm of n1, anchor=west] {Level 1 (1 true)};
\node[lbl, right=5.8cm of n6, anchor=west] {Level 2 (2 true)};
\node[lbl, right=5.8cm of nE, anchor=west] {Level 3 (3 true)};
\node[lbl, right=5.8cm of nF, anchor=west] {Level 4: Tautology};

% ── Legend ────────────────────────────────────────────────────
\node[draw=RSVPdeep, fill=RSVPlight, rounded corners=3pt,
      inner sep=5pt, font=\scriptsize,
      below right=0.5cm and -1cm of nF, anchor=north] {
  \begin{tabular}{ll}
    \color{RSVPaccent}\textbf{——} & NOR path: $\bot\to\mathrm{NOR}\to\neg A\to\mathrm{NAND}\to\top$\\
    \color{RSVPpurple}\textbf{——} & AND path: $\bot\to\mathrm{AND}\to\neg B\to\mathrm{OR}\to\top$\\
    \color{RSVPdeep}\textbullet\ light & Named connectives
  \end{tabular}
};

  \end{tikzpicture}
  \caption{The Hasse diagram of $B_2$: all 16 binary Boolean connectives
    ordered by truth-set inclusion. Bottom: contradiction ($\bot = 0000$).
    Top: tautology ($\top = 1111$). Highlighted amber path traces
    $\bot \to \mathrm{NOR} \to \neg A \to \mathrm{NAND} \to \top$.
    Highlighted purple path traces
    $\bot \to \mathrm{AND} \to \neg B \to \mathrm{OR} \to \top$.}
  \label{fig:hasse}
\end{figure}

\section{Logic as Accessibility Geometry}

The Boolean lattice $B_2$ has a natural accessibility interpretation.
Each connective $f \in B_2$ has an accessibility entropy:
\[
  S(f) = \log|R_f| = \log|\{x \in \{TT,TF,FT,FF\} : f(x) = 1\}|
\]
where $R_f$ is the truth region of $f$.

\begin{itemize}
  \item $S(\bot) = \log 0 = -\infty$ (contradiction: no true inputs)
  \item $S(\mathrm{AND}) = \log 1 = 0$ (most restrictive non-contradiction)
  \item $S(\mathrm{XOR}) = \log 2$ (two true inputs)
  \item $S(\mathrm{OR}) = \log 3$ (three true inputs)
  \item $S(\top) = \log 4$ (tautology: maximum accessibility)
\end{itemize}

The Hasse diagram is the accessibility landscape of the binary connective
space. Moving upward in the diagram corresponds to increasing accessibility
entropy. The accessibility gradient flow — the flow toward higher $S$ —
runs from contradiction to tautology along the edges of the lattice.

\begin{proposition}{Implication as Accessibility Ordering}{impl-access}
  $p \to q$ (logical implication) holds if and only if
  $S(p \wedge \cdot) \leq S(q \wedge \cdot)$ for all connectives:
  $p$'s truth region is a subset of $q$'s. Implication is
  accessibility containment.
\end{proposition}

\section{The Penteract and Hypercube Splitting}

The binary connective space $B_2 = Q_4$ can be extended to three
variables: $B_3 = Q_8$ with $2^8 = 256$ connectives. Visualizing $Q_8$
directly is impossible, but it decomposes into two copies of $Q_4$:

\[
  Q_5 = Q_4 \times \{0,1\}
\]

This is the \emph{penteract}: five-dimensional hypercube, visualized as
two tesseracts ($Q_4$) connected by matching edges. Each splitting
corresponds to fixing one binary variable and examining the resulting
four-variable space.

There are five coordinate splittings of $Q_5$ (one per dimension),
each revealing a different organizational principle of the space.
No single visualization captures all five simultaneously — any projection
into 2D discards information. This is the Boolean version of the
information loss theorem (Appendix~B): the geometry of logic cannot
be fully represented in any flat diagram.

\section{Seven Curricula, One Landscape}

Returning to the seven mathematical schools: each school produces a
different entry point into the Boolean landscape.

\begin{itemize}
  \item The \textbf{geometry-first} student sees the Hasse diagram as a
        polytope — a projection of a hypercube — and immediately asks
        what happens in higher dimensions.
  \item The \textbf{algebra-first} student sees a lattice with two
        operations ($\wedge$, $\vee$) and a complement, and asks about
        the algebraic axioms.
  \item The \textbf{logic-first} student sees a space of connectives
        ordered by implication and asks which paths correspond to
        valid inference chains.
  \item The \textbf{statistics-first} student sees each connective as
        a binary classifier and asks which achieves highest accuracy
        on which data distributions.
  \item The \textbf{calculus-first} student sees the lattice as a
        discrete gradient system and asks what the continuous limit looks like.
\end{itemize}

None of these perspectives is wrong. All of them reveal genuine structure
in the object. The logical geometry of $B_2$ is rich enough to contain
all five entry points simultaneously.

This is the key pedagogical claim: starting from any of the seven schools,
a sufficiently curious student will eventually discover the others.
The accessibility landscape of mathematics is connected — every concept
can be reached from every starting point — but the paths differ, and
different paths develop different cognitive orientations.

\begin{namedthm}{The Curriculum Equivalence Theorem}
  All seven mathematical schools are connected in the accessibility
  graph of mathematical concepts: from any starting concept, any other
  concept can be reached by a finite chain of natural conceptual
  extensions. The seven schools differ not in their ultimate destinations
  but in the order of traversal and the depth of each concept encountered
  along the way.
\end{namedthm}

\begin{exercise}
  Compute the Hamming distance between all pairs of named connectives
  ($\top, \bot, \mathrm{AND}, \mathrm{OR}, \mathrm{NAND}, \mathrm{NOR},
  \mathrm{XOR}, \mathrm{XNOR}, \neg A, \neg B$). Which pair is most
  logically similar (minimum distance)? Which is most dissimilar
  (maximum distance)?
\end{exercise}

\begin{exercise}
  The Boolean lattice $B_2$ has 16 elements. How many elements does
  $B_3$ (three-variable connectives) have? How many edges does the
  Hasse diagram of $B_3$ have? At what rate does the complexity of
  the Boolean landscape grow with the number of variables?
\end{exercise}

\begin{exercise}[title={Pedagogical}]
  Design a three-week logic-first mathematics curriculum for 10-year-olds.
  What would be taught in week 1 (the Boolean lattice), week 2
  (circuits and computation), and week 3 (connections to arithmetic)?
  What algebraic concepts would students encounter naturally, without
  being explicitly taught algebra?
\end{exercise}
